% 原书分章 PDF 第 2—6 页：章首及第 22.1 节；第 1 页部标题由总书管理。
\chapter{最优动态处理规则}
\label{ch:optimal-dynamic-treatment-rules}
\chapterauthors{Alexander R. Luedtke 与 Mark J. van der Laan}
\chapteraffiliations{%
A. R. Luedtke（通讯作者）：美国 Fred Hutchinson 癌症研究中心疫苗与传染病部，1100 Fairview Ave, Seattle, WA 98109, USA；电子邮箱：\href{mailto:aluedtke@fredhutch.org}{aluedtke@fredhutch.org}。\\
M. J. van der Laan：美国加利福尼亚大学伯克利分校生物统计学部及统计学系，101 Haviland Hall, \#7358, Berkeley, CA 94720, USA；电子邮箱：\href{mailto:laan@berkeley.edu}{laan@berkeley.edu}。\\
本章 DOI：\url{https://doi.org/10.1007/978-3-319-65304-4_22}。}

假定我们观测到某个随时间变化的随机变量的 $n$ 个独立同分布观测。这个随机变量依次包含基线协变量、初始处理与删失指示变量、中间协变量、后续处理与删失指示变量，以及最终结局。例如，这些数据可以来自一项序贯随机对照试验（sequential RCT）：研究者随访一组受试者，在两个时间点依次随机分配处理；一线处理的接受概率可以由某个基线协变量决定，而二线处理的接受概率可以由随时间变化的中间协变量（如某个感兴趣的生物标志物）决定。这类试验通常称为序贯多重分配随机试验（sequential multiple assignment randomized trial，SMART）。动态处理规则（dynamic treatment rule）把可获得的历史作为自变量，以确定性的方式分配处理。如果在两个时间点分配处理，那么动态处理规则就包含两个规则，每个时间点各一个。动态处理下的平均结局是我们关心的反事实量，表示如果每个人都按该动态处理规则接受处理，平均结局将会是多少。最优动态处理规则（optimal dynamic treatment rule）定义为使平均结局最大的动态处理规则。

本章末尾介绍的既有方法依赖半参数模型，对数据生成过程作出了很强的假设。我们则把数据分布的统计模型定义为非参数模型，仅可能利用有关处理机制（如随机对照试验中已知的处理机制）和删失机制的知识。为了不只考虑要求最高的完全最优规则，我们定义 $V$ 最优规则（$V$-optimal rule）：它是在可用协变量中，仅使用用户给定子集 $V$ 的最优规则。这样，我们就能考虑更容易估计的次优规则，进而对这些次优规则下的反事实平均结局进行统计推断。

\begin{sourcebox}
本章介绍两个时间点的 $V$ 最优规则下平均结局的半参数推断。我们将证明，最优规则下的平均结局是数据分布的路径可微参数，表明在一定条件下可构造其渐近线性估计量。事实上，我们得到一个令人意外的结果：该目标参数的路径导数，等于把给定动态处理规则设为最优规则、并视后者为已知时，反事实平均结局的路径导数。借助前面关于给定规则下平均结局的双重稳健且有效估计的结果（见第 4 章），即可得到最优规则下平均结局的交叉验证靶向最小损失估计（cross-validated targeted minimum loss-based estimation，CV-TMLE）。随后证明其渐近线性与有效性，从而为最优动态处理下的平均结局，或其相对于标准处理的对比，构造置信区间。
\end{sourcebox}

在 SMART 中，要得到渐近有效的统计推断，只需最优动态处理的估计量与最优动态处理本身之间的某个二阶差异渐近可忽略。如果只考虑依赖某个一维时变协变量的规则，或愿意作出光滑性假设，那么这一条件是合理的。估计最优平均结局时，这一条件似乎是必要的；但如果把目标参数重新定义为我们交叉验证得到的最优动态处理估计之下的平均结局的平均值，就不再需要这一条件。这个参数通过最优动态处理的估计而依赖数据，因此我们把它称为数据自适应参数（data-adaptive parameter）。

\section{最优动态处理估计问题}

为便于说明，本章只考虑两个时间点的处理。假定我们观测到下列随机变量的 $n$ 个独立同分布副本 $O_1,\ldots,O_n\in\mathcal O$：
\[
 O=(L(0),A(0),L(1),A(1),Y)\sim P_0,
\]
其中 $A(j)=(A_1(j),A_2(j))$，$A_1(j)$ 是二元处理，$A_2(j)$ 是在“时间”$j$ 尚未被右删失的指示变量，$j=0,1$。也就是说，$A_2(0)=0$ 意味着无法观测 $(L(1),A_1(1),Y)$，而 $A_2(1)=0$ 意味着无法观测 $Y$。每个时间点 $j$ 都有在处理之前测量的协变量 $L(j)$，$j=0,1$；感兴趣的结局记为 $Y$，发生在时间点 1 之后。对于随时间变化的过程 $X(\cdot)$，我们采用记号 $\bar X(t)=(X(s):s\leq t)$，并令 $\bar X(-1)=\varnothing$。令 $\mathcal M$ 为一个统计模型，它对 $L(0)$ 的边际分布 $Q_{0,L(0)}$ 以及给定 $A(0),L(0)$ 时 $L(1)$ 的条件分布 $Q_{0,L(1)}$ 不作假设，但可以对给定 $\bar A(j-1),\bar L(j)$ 时 $A(j)$ 的条件分布 $g_{0A(j)}$ 作出假设，$j=0,1$。我们把 $g_0$ 称为干预机制，它可以按如下方式分解为处理机制 $g_{01}$ 和删失机制 $g_{02}$：
\[
\begin{aligned}
 g_0(O)=\prod_{j=1}^{2}
 &g_{01}\bigl(A_1(j)\mid\bar A(j-1),\bar L(j)\bigr)\\
 &\quad g_{02}\bigl(A_2(j)\mid A_1(j),\bar A(j-1),\bar L(j)\bigr).
\end{aligned}
\]
特别地，数据可能来自 SMART；此时 $g_{01}$ 已知。

令 $V(1)$ 为 $(L(0),A(0),L(1))$ 的函数，$V(0)$ 为 $L(0)$ 的函数，并记 $V=(V(0),V(1))$。考虑分别用于分配处理 $A(0)$ 和 $A(1)$ 的动态处理规则
$V(0)\mapsto d_{A(0)}(V(0))\in\{0,1\}\times\{1\}$ 和
$(A(0),V(1))\mapsto d_{A(1)}(A(0),V(1))\in\{0,1\}\times\{1\}$。
其中，$A(0)$ 的规则仅是 $V(0)$ 的函数，$A(1)$ 的规则仅是 $(A(0),V(1))$ 的函数。注意，这些规则都要求把删失指示变量设为 $A_2(j)=1$，$j=0,1$。令 $\mathcal D$ 为所有这类规则组成的集合。我们假定 $V(0)$ 是 $V(1)$ 的函数（即观测到 $V(1)$ 就包含了观测到 $V(0)$），但将在下面的定理中给出一个替代假设。对 $d\in\mathcal D$，记
\[
 d(a(0),v)\equiv\bigl(d_{A(0)}(v(0)),d_{A(1)}(a(0),v(1))\bigr).
\]
如果假定这些变量满足如下结构方程模型（Pearl，2009a）：
\[
\begin{aligned}
 L(0)&=f_{L(0)}(U_{L(0)}),\\
 A(0)&=f_{A(0)}(L(0),U_{A(0)}),\\
 L(1)&=f_{L(1)}(L(0),A(0),U_{L(1)}),\\
 A(1)&=f_{A(1)}(\bar L(1),A(0),U_{A(1)}),\\
 Y&=f_Y(\bar L(1),\bar A(1),U_Y),
\end{aligned}
\]
其中函数集合 $f=(f_{L(0)},f_{A(0)},f_{L(1)},f_{A(1)})$ 未作指定或只作了部分指定，那么我们可以在修正后的系统中定义反事实 $Y_d$：把 $A(0),A(1)$ 的方程分别替换为 $A(0)=d_{A(0)}(V(0))$ 和 $A(1)=d_{A(1)}(A(0),V(1))$。将这些反事实量的分布记为 $P_{0,d}$；注意，$P_{0,d}$ 由函数集合 $f$ 与外生变量 $(U_{L(0)},U_{A(0)},U_{L(1)},U_{A(1)},U_Y)$ 的联合分布共同确定。现在可把 $P_{0,d}$ 下的因果最优规则定义为
$d_0^*=\argmax_{d\in\mathcal D}\E_{P_{0,d}}Y_d$。
如果假定序贯随机化成立，即给定 $L(0)$ 时 $A(0)$ 与 $U_{L(1)},U_Y$ 独立，且给定 $\bar L(1),A(0)$ 时 $A(1)$ 与 $U_Y$ 独立，那么可以利用 $g$ 计算公式，通过分布 $P_0$ 下的观测数据识别 $P_{0,d}$：
\begin{equation}
\begin{aligned}
 &p_{0,d}(L(0),A(0),L(1),A(1),Y)\\
 &\quad\equiv I\bigl(A=d(A(0),V)\bigr)q_{0,L(0)}(L(0))
 q_{0,L(1)}(L(1)\mid L(0),A(0))\\
 &\qquad\quad q_{0,Y}(Y\mid\bar L(1),\bar A(1)).
\end{aligned}
\tag{22.1}
\end{equation}
这里，$p_{0,d}$ 是 $P_{0,d}$ 的密度；$q_{0,L(0)}$、$q_{0,L(1)}$、$q_{0,Y}$ 分别是 $Q_{0,L(0)}$、$Q_{0,L(1)}$、$Q_{0,Y}$ 的密度，而 $Q_{0,Y}$ 表示给定 $\bar L(1),\bar A(1)$ 时 $Y$ 的分布。我们假定以上所有密度都相对于某个控制测度 $\mu$ 绝对连续。在 Neyman--Rubin 因果模型下，也有类似的可识别性结果和 $g$ 计算公式（Robins，1987）。更一般地，对分布 $P\in\mathcal M$，可把 $g$ 计算分布 $P_d$ 定义为具有如下密度的分布：
\[
\begin{aligned}
 &p_d(L(0),A(0),L(1),A(1),Y)\\
 &\quad\equiv I\bigl(A=d(A(0),V)\bigr)q_{L(0)}(L(0))
 q_{L(1)}(L(1)\mid L(0),A(0))\\
 &\qquad\quad q_Y(Y\mid\bar L(1),\bar A(1)),
\end{aligned}
\]
其中 $q_{L(0)}$、$q_{L(1)}$ 和 $q_Y$ 是 $P$ 下与 $q_{0,L(0)}$、$q_{0,L(1)}$ 和 $q_{0,Y}$ 对应的量。

在本章后续部分，对于静态或动态干预 $d$，使用记号 $L_d$（或 $Y_d,O_d$）时，均指概率分布为式 (22.1) 中 $P_d$ 的随机变量，因此我们讨论的所有量都是统计参数。例如，下一定理中定义的 $\E_0(Y_{a(0)a(1)}\mid V_{a(0)}(1))$，表示在概率分布 $P_{0,a(0)a(1)}$ 下，给定 $V_{a(0)}(1)$ 时 $Y_{a(0)a(1)}$ 的条件期望；这个分布就是以上针对静态干预 $(a(0),a(1))$ 给出的 $g$ 计算分布。此外，只要写出某个 $P_d$ 下的这些参数，我们就自动假定 $P$ 处满足使 $g$ 计算公式有良好定义所需的正值性假设。为此，假定 $P$ 处满足下列正值性条件就已足够：
\begin{equation}
\begin{aligned}
 \Pr_P\left(0<\min_{a_1\in\{0,1\}}g_{0A(0)}(a_1,1\mid L(0))\right)&=1,\\
 \Pr_P\left(0<\min_{a_1\in\{0,1\}}g_{0A(1)}(a_1,1\mid\bar L(1),A(0))\right)&=1.
\end{aligned}
\tag{22.2}
\end{equation}
强正值性假设定义为上述假设中把 0 替换为某个 $\delta>0$ 所得的条件。

现在定义一个统计参数，表示 $P_d$ 下结局 $Y_d$ 的均值。对任意规则 $d\in\mathcal D$，令
\[
 \Psi_d(P)\equiv\E_{P_d}Y_d.
\]
对于分布 $P$，把 $V$ 最优规则定义为
\[
 d_P=\argmax_{d\in\mathcal D}\E_{P_d}Y_d.
\]
为简便起见，把 $P_0$ 下的 $V$ 最优规则 $d_{P_0}$ 写为 $d_0$。定义参数映射 $\Psi:\mathcal M\to\mathbb R$ 为 $\Psi(P)=\E_{P_{d_P}}Y_{d_P}$。本章第一部分研究下列参数的推断：
\[
 \psi_0\equiv\Psi(P_0)=\E_{P_{0,d_0}}Y_{d_0}.
\]
在我们的可识别性假设下，$d_0$ 等于因果最优规则 $d_0^*$。即使序贯随机化假设不成立，统计参数 $\psi_0$ 本身仍是一个有意义的研究对象。在后文中，我们不再讨论序贯随机化假设。

下一个定理将 $V$ 最优个体化处理规则 $d_0$ 表示为 $P_0$ 的显式函数。

\begin{bookresult}{定理}{22.1}
假定 $V(0)$ 是 $V(1)$ 的函数。$V$ 最优规则 $d_0$ 可表示为 $P_0$ 的下列显式参数：
\[
\begin{aligned}
 \bar Q_{b,20}(a(0),v(1))
 &=\E_0(Y_{a(0),A(1)=(1,1)}\mid V_{a(0)}(1)=v(1))\\
 &\quad-\E_0(Y_{a(0),A(1)=(0,1)}\mid V_{a(0)}(1)=v(1)),\\
 d_{0,A(1)}(A(0),V(1))
 &=\bigl(I(\bar Q_{b,20}(A(0),V(1))>0),1\bigr),\\
 \bar Q_{b,10}(v(0))
 &=\E_0(Y_{(1,1),d_{0,A(1)}}\mid V(0))
 -\E_0(Y_{(0,1),d_{0,A(1)}}\mid V(0)),\\
 d_{0,A(0)}(V(0))
 &=\bigl(I(\bar Q_{b,10}(V(0))>0),1\bigr),
\end{aligned}
\]
其中 $a(0)\in\{0,1\}\times\{1\}$。如果 $V(1)$ 不包含 $V(0)$，但对所有
$(a(0),a(1))\in\{\{0,1\}\times\{1\}\}^2$，都有
\begin{equation}
 \E_0(Y_{a(0),a(1)}\mid V(0),V_{a(0)}(1))
 =\E_0(Y_{a(0),a(1)}\mid V_{a(0)}(1)),
\tag{22.3}
\end{equation}
那么上述 $V$ 最优规则 $d_0$ 的表达式仍然成立。
\end{bookresult}

沿用 Robins（2004）的术语，我们把 $\bar Q_{b,10}$ 和 $\bar Q_{b,20}$ 称为（第一和第二时间点的）处理增益函数（blip function）。
